\begin{table}[!htbp]
\centering
\caption{Close-Election Event Study of Moody's Rating Levels (Pre-Trend Diagnostic)}
\label{tab:post_northeast_rdd_event_study}
\small
\begin{tabular}{lc}
\toprule
 & Rating gap (notches) \\
\midrule
$h=-2$ & -0.076 \\
 & (0.139) \\
$h=-1$ & 0 (reference) \\
$h=0$ & 0.126 \\
 & (0.137) \\
$h=1$ & -0.070 \\
 & (0.120) \\
$h=2$ & 0.165 \\
 & (0.158) \\
\midrule
Observations & 858 \\
Events & 220 \\
Calendar-year fixed effects & Yes \\
\bottomrule
\end{tabular}
\vspace{0.05in}
\begin{minipage}{0.96\textwidth}\footnotesize\emph{Notes:} Two-group event study of the Moody's rating level (notches) on event-time indicators interacted with \emph{close success}, the common event-time path, and the close-success level, with calendar-year fixed effects and standard errors clustered by municipality. The reference period is $h=-1$ (the election fiscal year). Positive values mean close successes hold higher ratings than close failures. The $h=-2$ coefficient is the pre-trend test. * p $<$ 0.10, ** p $<$ 0.05, *** p $<$ 0.01.\end{minipage}
\end{table}
